\documentclass[11pt]{article}
\newcommand{\DocShort}{Gu\'ia de lectura}
\usepackage{fontspec}
\usepackage[spanish,es-noshorthands]{babel}
\decimalpoint
\usepackage{geometry}
\geometry{letterpaper, top=2.5cm, bottom=2.5cm, left=2.8cm, right=2.8cm}
\usepackage[expansion=false]{microtype}
\usepackage{parskip}
\setlength{\parskip}{6pt}
\usepackage{amsmath,amssymb}
\usepackage{booktabs}
\usepackage{array}
\usepackage{longtable}
\usepackage{caption}
\usepackage[dvipsnames,table]{xcolor}
\usepackage{enumitem}
\setlist[itemize]{topsep=2pt, itemsep=2pt, parsep=0pt}
\setlist[enumerate]{topsep=2pt, itemsep=2pt, parsep=0pt}
\usepackage{fancyhdr}
\usepackage[hidelinks,pdfencoding=auto]{hyperref}
\usepackage{titlesec}
\usepackage{tcolorbox}
\tcbuselibrary{breakable}
\usepackage{pdfpages}

\definecolor{darkblue}{HTML}{1B3A5C}
\definecolor{medblue}{HTML}{2E6DA4}
\definecolor{lightblue}{HTML}{D6E8F7}
\definecolor{teal}{HTML}{1A7A6E}
\definecolor{lightteal}{HTML}{D0EDEA}
\definecolor{lightgrey}{HTML}{F5F5F5}
\definecolor{midgrey}{HTML}{6F6F6F}
\definecolor{warnyellow}{HTML}{FFF3CD}

\titleformat{\section}
{\large\bfseries\color{darkblue}}
{\thesection}{0.75em}{}[{\color{medblue}\titlerule[0.8pt]}]
\titleformat{\subsection}
{\normalsize\bfseries\color{darkblue}}
{\thesubsection}{0.75em}{}
\titleformat{\subsubsection}
{\normalsize\bfseries\color{medblue}}
{\thesubsubsection}{0.75em}{}

\pagestyle{fancy}
\fancyhf{}
\fancyhead[L]{\small\color{midgrey}ICV 442 Acueductos y Alcantarillados}
\fancyhead[R]{\small\color{midgrey}\DocShort}
\fancyfoot[C]{\small\color{midgrey}\thepage}
\renewcommand{\headrulewidth}{0.4pt}
\renewcommand{\footrulewidth}{0pt}

\rowcolors{2}{lightgrey}{white}
\renewcommand{\arraystretch}{1.18}

\newtcolorbox{infobox}{
  breakable,
  colback=lightblue,
  colframe=medblue,
  boxrule=0.8pt,
  arc=3pt,
  left=8pt,
  right=8pt,
  top=7pt,
  bottom=7pt
}

\newtcolorbox{warnbox}{
  breakable,
  colback=warnyellow,
  colframe=orange!70!black,
  boxrule=0.8pt,
  arc=3pt,
  left=8pt,
  right=8pt,
  top=7pt,
  bottom=7pt
}

\newcommand{\doccover}[3]{
  \begin{center}
  {\small\color{midgrey}Pontificia Universidad Cat\'olica Madre y Maestra\\Escuela de Ingenier\'ia Civil y Ambiental}

  \vspace{1.0cm}
  {\color{medblue}\rule{\linewidth}{2pt}}
  \vspace{0.3cm}

  {\Large\bfseries\color{darkblue}ICV-442-T Acueductos y Alcantarillado}\\[0.35em]
  {\LARGE\bfseries\color{darkblue}#1}\\[0.25em]
  {\large\itshape #2}

  \vspace{0.3cm}
  {\color{medblue}\rule{\linewidth}{2pt}}
  \vspace{0.9cm}
  \end{center}

  \begin{center}
  \begin{tabular}{ll}
  \toprule
  \textbf{Asignatura} & ICV-442-T Acueductos y Alcantarillado \\
  \textbf{Profesor} & Ricardo Rafael Hern\'andez Moreira, PhD \\
  \textbf{Periodo acad\'emico} & Mayo--agosto de 2026 \\
  \textbf{Proyecto} & Dise\~no de sistemas de saneamiento para Las Charcas, Azua \\
  \textbf{Versi\'on} & #3 \\
  \bottomrule
  \end{tabular}
  \end{center}

  \vspace{0.4cm}
}


\begin{document}
\doccover{Gu\'ia de lectura del proyecto}{Ruta de uso del paquete documental}{18 de mayo de 2026}

\begin{infobox}
Este documento indica qu\'e leer primero, para qu\'e sirve cada archivo y cu\'ales son las fechas y productos que estructuran el proyecto. El paquete est\'a pensado para que el estudiantado trabaje con documentos separados y no tenga que navegar un solo archivo extenso.
\end{infobox}

\section{Orden recomendado de lectura}
\begin{enumerate}
\item \textbf{T\'erminos de referencia}. Define el encargo, el alcance, los entregables, la organizaci\'on del trabajo y la convenci\'on de nombres.
\item \textbf{Especificaciones t\'ecnicas de dise\~no}. Establece los criterios m\'inimos para agua potable, alcantarillado sanitario y drenaje pluvial.
\item \textbf{Cronograma del proyecto}. Traduce el encargo en hitos semanales, entregas y consecuencias por retraso.
\item \textbf{R\'ubrica de evaluaci\'on}. Indica c\'omo se calificar\'a el proyecto escrito y la presentaci\'on oral.
\item \textbf{Lista de cotejo}. Se usa al final para verificar que la entrega est\'a completa y bien organizada.
\item \textbf{Anexos operativos}. Incluyen patrones para EPANET y SWMM, as\'i como la plantilla de acta y bit\'acora.
\end{enumerate}

\section{Archivos del paquete}
\begin{center}
\begin{tabular}{p{0.18\linewidth} p{0.31\linewidth} p{0.39\linewidth}}
\toprule
\textbf{Archivo} & \textbf{Uso principal} & \textbf{Cu\'ando usarlo} \\
\midrule
01\_Guia\_de\_lectura\_del\_proyecto.pdf & Orientaci\'on general & Antes de empezar y cuando se necesite ubicar un tema \\
02\_Terminos\_de\_referencia.pdf & Reglas del proyecto & Durante la planificaci\'on inicial y al organizar entregas \\
03\_Especificaciones\_tecnicas\_de\_diseno.pdf & Criterios de dise\~no & Durante c\'alculos, modelaci\'on y redacci\'on t\'ecnica \\
04\_Cronograma\_del\_proyecto.pdf & Hitos y fechas & Toda la duraci\'on del proyecto \\
05\_Rubrica\_de\_evaluacion.pdf & Criterios de calificaci\'on & Antes de cada entrega parcial y antes de la defensa oral \\
06\_Lista\_de\_cotejo\_entrega\_final.pdf & Verificaci\'on final & En la semana de entrega final \\
07\_Anexo\_patron\_demanda\_EPANET.pdf & Patr\'on horario de demanda & Al configurar el modelo de agua potable \\
08\_Anexo\_patron\_descarga\_SWMM.pdf & Patr\'on horario sanitario & Al configurar el modelo sanitario \\
09\_Anexo\_acta\_y\_bitacora.pdf & Formatos de gesti\'on & En la conformaci\'on del equipo y el seguimiento semanal \\
12\_Formulario\_revision\_por\_pares.pdf & Retroalimentaci\'on entre equipos & En la semana de revisi\'on por pares y antes del cierre final \\
13\_Formato\_presentacion\_semanal.pdf & Presentaci\'on de avance de 5 minutos & Antes de cada jueves de seguimiento y ensayo breve del equipo \\
\bottomrule
\end{tabular}
\end{center}

\section{Fechas y productos clave}
\begin{center}
\begin{tabular}{p{0.17\linewidth} p{0.28\linewidth} p{0.39\linewidth}}
\toprule
\textbf{Semana} & \textbf{Fecha de referencia} & \textbf{Producto} \\
\midrule
Semana 3 & Viernes 22 de mayo de 2026 & Acta de conformaci\'on del equipo \\
Semana 4 & Viernes 29 de mayo de 2026 & Memoria de proyecci\'on poblacional \\
Semana 5 & Viernes 5 de junio de 2026 & Tabla maestra de par\'ametros y brechas de informaci\'on \\
Semana 6 & Viernes 12 de junio de 2026 & Anteproyecto con trazado preliminar \\
Semana 7 & Viernes 19 de junio de 2026 & Reporte preliminar EPANET \\
Semana 8 & Viernes 26 de junio de 2026 & Reporte preliminar SWMM \\
Semana 9 & Viernes 3 de julio de 2026 & Borrador integrado \\
Semana 10 & Viernes 10 de julio de 2026 & Revisi\'on por pares, requisito habilitante \\
Semana 12 & Viernes 24 de julio de 2026, 23:59 & Entrega final del proyecto escrito \\
Semana 14 & 3 y 4 de agosto de 2026 & Presentaciones orales finales \\
\bottomrule
\end{tabular}
\end{center}

\section{Reglas pr\'acticas de uso}
\begin{itemize}
\item Ning\'un documento sustituye a otro. Los t\'erminos de referencia y las especificaciones deben leerse conjuntamente.
\item Las decisiones t\'ecnicas deben justificarse con normativa aplicable. La regla general es priorizar INAPA y usar MOPC cuando sea m\'as restrictivo.
\item Los archivos finales del equipo deben seguir la convenci\'on de nombres indicada en los t\'erminos de referencia.
\item Los anexos operativos no son lecturas introductorias; se consultan cuando el trabajo ya entr\'o en fase de modelaci\'on o seguimiento.
\end{itemize}

\section{Sugerencia de trabajo}
\begin{warnbox}
Se recomienda que cada equipo empiece por una lectura r\'apida del paquete completo y luego asigne responsables temporales para tres frentes: informaci\'on b\'asica y normativa, modelaci\'on hidr\'aulica, y documentaci\'on gr\'afica. Esa distribuci\'on no divide el proyecto por sistemas; solo ordena el trabajo.
\end{warnbox}

\end{document}
